\section{Week 18}
\sol{18.1}{Apply the one-step estimate with $n=0$, $1$ and $2$ in turn. Each time, substitute the bound already found, multiplied by $q$; since $q>0$, multiplying keeps the direction of the inequality:
\begin{align*}
 e_1&\le qe_0+\eps_0,\\
 e_2&\le qe_1+\eps_1\le q(qe_0+\eps_0)+\eps_1=q^2e_0+q\eps_0+\eps_1,\\
 e_3&\le qe_2+\eps_2\le q(q^2e_0+q\eps_0+\eps_1)+\eps_2=q^3e_0+q^2\eps_0+q\eps_1+\eps_2.
\end{align*}
In the bound for $e_n$ in part~1 of Theorem~\ref{thm:inexact}, the update error $\eps_j$ carries the weight $q^{n-1-j}$. So the earliest update error $\eps_0$ has weight $q^{n-1}$, and the latest one, $\eps_{n-1}$, has weight $q^0=1$.

These are the right powers because of how many updates follow each error. The error $\eps_j$ is made in the update that produces $y_{j+1}$. After it come the $n-1-j$ updates that produce $y_{j+2},\ldots,y_n$, and each of them multiplies its effect by $q$. The error $\eps_0$ is followed by $n-1$ updates, and the error $\eps_{n-1}$ by none. For $n=3$ this gives the weights $q^2$, $q$ and $1$ of the bound for $e_3$ above, a useful check on the exponent $n-1-j$.}

\sol{18.2}{\emph{The formula.} At $n=0$ the formula gives $\eta(1-q^0)/(1-q)=0=y_0$. Suppose that $y_n=\eta(1-q^n)/(1-q)$ for some $n$. Then
\[
 y_{n+1}=q\eta\,\frac{1-q^n}{1-q}+\eta
 =\eta\,\frac{q-q^{n+1}+1-q}{1-q}
 =\eta\,\frac{1-q^{n+1}}{1-q},
\]
which is the formula for $n+1$. By induction it holds for every $n\in\NN$.

\emph{Sharpness.} The map $T(x)=qx$ is a contraction of $\RR$ with factor $q$, its fixed point is $0$, and every update error is $|y_{n+1}-qy_n|=\eta$. So the example satisfies the hypotheses of part~2 of Theorem~\ref{thm:inexact} with this $\eta$. Its starting error is $e_0=0$. Since $\eta\ge0$ and $q^n\le1$, every $y_n$ is nonnegative, so $e_n=y_n=\eta(1-q^n)/(1-q)$.

Now let $\eta>0$ and let $K$ be any number smaller than $1/(1-q)$. If the factor $1/(1-q)$ in part~2 could be replaced by $K$, this example would satisfy $e_n\le q^ne_0+K\eta=K\eta$ for every $n\ge1$, that is,
\[
 \eta\,\frac{1-q^n}{1-q}\le K\eta.
\]
Dividing by $\eta>0$ and multiplying by $1-q>0$ turns this into $1-q^n\le K(1-q)$, or $q^n\ge1-K(1-q)$. Since $K(1-q)<1$, the number $c=1-K(1-q)$ is positive. Because $q^n\to0$, there is a threshold beyond which $q^n<c$, and for every such $n$ the required inequality fails. So the smaller factor $K$ fails for this example, at every step from some point on. No factor smaller than $1/(1-q)$ works for every iteration that satisfies the hypotheses, which is what it means for the factor $1/(1-q)$ to be sharp.}

\sol{18.3}{\emph{The hypotheses.} The real line with the distance $|x-y|$ is nonempty, and it is complete; the completeness of the real numbers is one of the starting facts of Chapter~\ref{ch:limits}. The map $T(x)=0$ sends every real number to the real number $0$, so it maps $\RR$ into itself. For every $q$ with $0<q<1$ and all real $x$ and $y$,
\[
 |T(x)-T(y)|=|0-0|=0\le q|x-y|,
\]
so $T$ is a contraction with factor $q$. Its fixed point is $0$, because $T(0)=0$. Every $y_n$ is a real number, so it lies in $X=\RR$. Finally,
\[
 |y_{n+1}-T(y_n)|=|(-1)^{n+1}\eta-0|=\eta,
\]
so the update errors are $\eps_n=\eta\ge0$, and the hypothesis of part~2 holds with this same number $\eta$.

\emph{The bound of part~2.} Every error is $e_n=|(-1)^n\eta-0|=\eta$. For $n\ge1$, the bound of part~2 is
\[
 q^n\eta+\eta\,\frac{1-q^n}{1-q}=q^n\eta+\eta\left(1+q+\cdots+q^{n-1}\right).
\]
The sum in parentheses contains the term $1$, and all its other terms and the term $q^n\eta$ are nonnegative, so the bound is at least $\eta=e_n$. The bound holds, as the theorem promises.

\emph{No convergence.} Consecutive terms are always $2\eta$ apart:
\[
 |y_{n+1}-y_n|=|(-1)^{n+1}\eta-(-1)^n\eta|=|{-2}(-1)^n\eta|=2\eta.
\]
Take the tolerance $\eta$, which is positive. For every proposed threshold $N$, the indices $N$ and $N+1$ are both at least $N$, and $|y_{N+1}-y_N|=2\eta\ge\eta$. So no threshold works, and the sequence is not a Cauchy sequence. By Proposition~\ref{ch11:prop:convergent-cauchy}, every convergent sequence is a Cauchy sequence, so this sequence does not converge.

\emph{Why the bound does not settle convergence.} The error bound controls one number at each step, the distance from $y_n$ to $x_*$. Here that distance is $\eta$ at every step, so the bound holds, and the distances even form a constant sequence. But two different points can lie at the same distance from $x_*$, here $\eta$ and $-\eta$, and the points may keep jumping between them. Knowing the distances to $x_*$ forces convergence only when those distances tend to zero, and with update errors that stay at $\eta$, the bound of part~2 does not tend to zero. That is why part~3 of the theorem needs the stronger hypothesis $\eps_n\to0$, which fails here.}

\sol{18.4}{\emph{The bound.} Every weight $q^{n-1-j}$ in part~1 of Theorem~\ref{thm:inexact} is positive, so we may replace each $\eps_j$ by the number $aq^j$, which is at least as large:
\begin{align*}
 \sum_{j=0}^{n-1}q^{n-1-j}\eps_j
 &\le\sum_{j=0}^{n-1}a q^{n-1-j}q^j\\
 &=\sum_{j=0}^{n-1}a q^{n-1}\\
 &=a n q^{n-1}.
\end{align*}
The exponent is $n-1-j+j=n-1$ in every term, and there are $n$ terms, one for each $j$ from $0$ to $n-1$. Part~1 therefore gives $e_n\le q^n e_0+a n q^{n-1}$ for every $n\ge1$.

\emph{The factor $n$.} We still have to show that $nq^{n-1}$ tends to zero. It is not enough that $q^{n-1}\to0$, because $q^{n-1}$ is multiplied by the growing number $n$. Put $r=(1+q)/2$. Since $q<1$, we have $q=\frac{q+q}2<\frac{1+q}2=r<\frac{1+1}2=1$, so $q<r<1$. Put $t=q/r$; then $0<t<1$, because $0<q<r$. Since $0<t<1$, the powers $1,t,\ldots,t^{n-1}$ decrease, so each of these $n$ numbers is at least the last one, $t^{n-1}$. Adding them up gives
\[
 n t^{n-1}\le1+t+\cdots+t^{n-1}=\frac{1-t^n}{1-t}\le\frac1{1-t}.
\]
Multiply by the positive number $r^{n-1}$ and use $q=rt$, so that $q^{n-1}=r^{n-1}t^{n-1}$:
\[
 0\le n q^{n-1}=r^{n-1}\,n t^{n-1}
 \le\frac{r^{n-1}}{1-t}.
\]
The number $1-t$ is fixed and positive, and $r^{n-1}$ becomes as small as we like, because $0<r<1$.

\emph{The whole bound tends to zero.} Write $b_n=q^ne_0+anq^{n-1}$ for the upper bound, and let $\delta>0$. We find a threshold beyond which each of the two terms of $b_n$ is less than $\delta/2$.
\begin{itemize}
\item If $e_0=0$, the first term is zero; put $N_1=1$. Otherwise, $q^n\to0$ (Section~\ref{ch11:sec:geometric}) gives a threshold $N_1$ with $q^n<\delta/(2e_0)$, and so $q^ne_0<\delta/2$, for every $n\ge N_1$.
\item If $a=0$, the second term is zero; put $N_2=1$. Otherwise, $r^m\to0$ gives a threshold $M$ with $r^m<(1-t)\delta/(2a)$ for every $m\ge M$. Put $N_2=M+1$. For every $n\ge N_2$ we have $n-1\ge M$, and so $anq^{n-1}\le ar^{n-1}/(1-t)<\delta/2$.
\end{itemize}
For every $n$ that is at least both $N_1$ and $N_2$, we get $0\le b_n<\delta/2+\delta/2=\delta$. So $b_n\to0$.

Since $0\le e_n\le b_n$, also $e_n\to0$, that is, $y_n\to x_*$. This conclusion also follows from part~3 of the theorem, because $aq^j\to0$. What the bound $b_n$ adds is information about speed: it says how small the error is after a given number of steps.}

\sol{18.5}{Several students may now receive the same task, as long as its capacity is not exceeded, so an assignment with capacities is not a matching in the original graph. We translate the problem into one about ordinary matchings by giving every unit of capacity its own vertex, called a \emph{slot}.

\emph{The condition is necessary.} Suppose an assignment with capacities exists, and let $S$ be any set of students. Every student in $S$ receives a task that they accept, and that task belongs to $N(S)$. For each task $r\in N(S)$, let $m_r$ be the number of students in $S$ who receive $r$; then $m_r\le c_r$. Every student in $S$ is counted in exactly one of the numbers $m_r$, so
\[
 |S|=\sum_{r\in N(S)}m_r\le\sum_{r\in N(S)}c_r.
\]
This holds for every set $S$ of students, not only for the whole class.

\emph{The condition is sufficient: a new graph.} Now suppose the inequality holds for every $S\subseteq L$. Build a new bipartite graph. Its students are those of $L$. In place of each task $r$ it has $c_r$ different slots $(r,1),(r,2),\ldots,(r,c_r)$; a task of capacity zero gets no slots at all. A student $s$ is joined to the slot $(r,i)$ exactly when $s$ accepts the task $r$. For example, if task~$1$ has capacity $2$ and task~$2$ has capacity $1$, the slots are $(1,1)$, $(1,2)$ and $(2,1)$, and a student who accepts task~$1$ is joined to both $(1,1)$ and $(1,2)$. The new graph is finite, because $L$ and $R$ are finite and each capacity is a whole number.

\emph{Hall's condition in the new graph.} Let $S\subseteq L$. A slot $(r,i)$ is joined to some student of $S$ exactly when some student of $S$ accepts $r$, that is, exactly when $r\in N(S)$. So the neighbors of $S$ in the new graph are all the slots of all the tasks in $N(S)$. Slots of different tasks are different, because their first coordinates differ, and task $r$ has exactly $c_r$ slots. Hence $S$ has exactly $\sum_{r\in N(S)}c_r$ neighbors in the new graph, and by assumption this number is at least $|S|$. This is Hall's condition~\eqref{eq:hall} for the new graph.

\emph{Back to tasks.} By Hall's theorem (Theorem~\ref{thm:hall}), the new graph has a matching that saturates $L$: every student receives a slot they are joined to, and no slot goes to two students. Give each student the task $r$ of the slot $(r,i)$ they received. The student accepts $r$, because they are joined to that slot. Task $r$ is given to at most $c_r$ students, because different students received different slots and $r$ has only $c_r$ slots. So this is an assignment with capacities. The translation preserves both what the students accept and the capacities, which is why solving the problem in the new graph solves the original one.

When every capacity equals $1$, each task has exactly one slot, the new graph is just the original graph, and the condition reads $|S|\le|N(S)|$. This is Hall's theorem again.}

\sol{18.6}{Here is one possible record.

\emph{Question.} The map $T(x)=(3x+1)/4$ is a contraction of $\RR$ with factor $q=3/4$, and its fixed point is $1$. If every update is computed with an error of at most $1/100$, do the points still converge to $1$? If not, how far from $1$ can they be?

\emph{Background.} Arithmetic with fractions; induction (Chapter~\ref{ch:induction}); distance, the triangle inequality, convergence, and the finite geometric sum (Chapter~\ref{ch:limits}); the contraction theorem (Chapter~\ref{ch:contraction}).

\emph{Outline of the proof.} (1) Insert the exact update $T(y_n)$ between $y_{n+1}$ and the fixed point, and use the triangle inequality, the equation $T(1)=1$ and the contraction inequality to get $e_{n+1}\le qe_n+\eps_n$. (2) Unfold this by induction to $e_n\le q^ne_0+\sum_{j=0}^{n-1}q^{n-1-j}\eps_j$. (3) Bound every update error by $\eta$ and add up the geometric weights, which gives $e_n\le q^ne_0+\eta(1-q^n)/(1-q)$. Here this reads $|y_n-1|\le(3/4)^n+\frac1{25}\bigl(1-(3/4)^n\bigr)$.

\emph{Counterexample to the first guess.} With the constant bias $1/100$, the points are $y_n=\frac{26}{25}\bigl(1-(3/4)^n\bigr)$, a formula proved by induction in Section~\ref{ch18:sec:experiment}. They converge to $26/25$, not to the fixed point $1$. So small update errors do not guarantee convergence to the fixed point.

\emph{An exact check.} At $n=1$ the iterate is $y_1=\frac{3\cdot0+1}4+\frac1{100}=\frac{26}{100}=\frac{13}{50}$, in agreement with the closed form $\frac{26}{25}\bigl(1-\frac34\bigr)=\frac{26}{100}$. The true error is $\bigl|\frac{13}{50}-1\bigr|=\frac{37}{50}$, and the certificate is $\frac34+\frac1{25}\cdot\frac14=\frac34+\frac1{100}=\frac{19}{25}=\frac{38}{50}$. So $\frac{37}{50}\le\frac{38}{50}$, computed exactly. This check confirms the formulas at one index only. The closed form and the certificate hold for every $n$ because they were proved.

\emph{What is new.} The project reconstructs and studies a standard bound for contractions whose updates are computed with errors, together with exact examples. It does not establish a new theorem compared with the published literature.

\emph{A further question.} Suppose $X=\RR$ and the update errors alternate in sign, so that $y_{n+1}-T(y_n)$ equals $\eta,-\eta,\eta,-\eta,\ldots$ in turn. After many steps, is the distance from the fixed point then bounded by a smaller multiple of $\eta$ than the $\eta/(1-q)$ allowed by part~2 of Theorem~\ref{thm:inexact}? For $T(x)=(3x+1)/4$ and $\eta=1/100$, Section~\ref{ch18:sec:experiment} reports that the points end up only $1/175$ away from $1$, far less than the $1/25$ that the certificate allows. The proofs in this chapter cannot answer the question, because the triangle inequality ignores the signs of the update errors and so cannot take advantage of their cancellation.}
